\section{实验、分析与结果}

% 本章是论文的核心，通常 5000–8000 字。
% 先展示结果，再进行分析和讨论。

\subsection{实验设计}

% 说明实验的具体设置、变量控制和实施过程。
% 如果是问卷研究，说明量表来源和信效度检验结果。

\subsection{结果展示}

% 用图表呈现主要结果。
% 图表标题使用中英双语，引用时使用 \ref{}：
% 如表~\ref{tab:example} 所示……
%
% 表格示例：
% \begin{table}[htbp]
%   \bicaption{描述性统计结果}{Descriptive Statistics}
%   \label{tab:example}
%   \centering
%   \begin{tabular}{lcc}
%     \topline
%     变量 & 均值 & 标准差 \\
%     \midrule
%     X1   & 3.45 & 0.82 \\
%     X2   & 4.12 & 0.67 \\
%     \bottomline
%   \end{tabular}
% \end{table}

\subsection{结果分析}

% 对结果进行解释，与文献中的发现进行对比。
% 指出哪些假设得到支持，哪些没有。

\subsection{讨论}

% 讨论结果的理论和实践意义。
% 分析可能的原因，承认局限性。
% 与第 2 章的文献综述形成呼应。
